\begin{textsample}{POS Dim 1 – vem\_ed\_24 – Score 5.00 – vem\_ed\_24/t125.txt}  \label{ex:f1_pos_057}
Leitores Osvaldo Succi Jr.
Com muita satisfação, realizamos o International Virtual Exchange Symposium (IVES Cesu 2024) na plataforma MS Teams e obtivemos um recorde de participantes: 246, \textbf{sendo} 43 estrangeiros, de 12 países das Américas, da Europa e da África.
A presença feminina deu o tom à \textbf{terceira} edição do IVES Cesu – tanto em palestrantes: seis de cinco diferentes países, \textbf{quanto} em ouvintes, \textbf{que} \textbf{representaram} dois terços da audiência.
O IVES Cesu 2024 \textbf{foi} diverso \textbf{quanto} a faixas etárias: 233 dos 246 participantes têm entre 26 e 65 anos, 6 têm de 18 a 25 anos e 7 estão acima de 65 anos.
Envolveu 111 professores de Fatec, 20 de Etec, 10 do Centro Paula Souza e 62 de instituições brasileiras externas ao CPS.
Esses dados mostram \textbf{que} devemos \textbf{seguir} \textbf{promovendo} eventos inclusivos e acessíveis a diferentes públicos.
E por falar em eventos, a agenda do segundo semestre de 2024 está repleta deles: Internacionalização em Casa \textbf{é} a pauta da Rede de Apoio ao Ensino Superior (RedAES), prevista para 29 de agosto; o 4º Congresso da Red LatAM COIL, on-line, vai de 4 a 6 de setembro e a principal conferência mundial de Intercâmbios Virtuais, IVEC, ocorre neste ano também de \textbf{forma} on-line, entre 21 e 24 de outubro.
Neste \textbf{número}, inauguramos a seção “Artigo de Opinião”, com \textbf{reflexões} de Priscilla Ferro sobre a trajetória do Teletandem aos PCIs.
Quer contribuir você também?
Confira a plataforma de submissões: https://publicacoescesu.cps.sp.gov.br/VEm/about/submissions Dúvidas sobre como colaborar?
Escreva para cesu.pci@fatec.sp.gov.br Boa leitura!

% matched lemmas: forma, número, promover, quanto, que, reflexão, representar, seguir, ser, terceiro
\end{textsample}
